
This section describes the experimental framework 
designed to reveal and interpret memorization in 
CNNs applied to image classification. The audit 
pipeline integrates model training, memorization 
analysis, adversarial evaluation, and explainability 
methods into a unified workflow.

\subsubsection{Dataset}

The CIFAR-10 dataset (Krizhevsky et al., 2009) is 
used as the primary benchmark for the experimental 
evaluation. CIFAR-10 consists of 60,000 color images 
of size $32 \times 32$ pixels, distributed across 
10 classes with 6,000 images per class. The dataset 
is partitioned into 50,000 training images and 10,000 
test images. To evaluate the model's behavior on 
out-of-distribution inputs, a set of custom images 
collected by the author is additionally used as a 
supplementary test set.

The input images are preprocessed as follows. Each 
image $x \in \mathbb{R}^{32 \times 32 \times 3}$ is 
normalized using the channel-wise mean 
$\mu = (0.4914, 0.4822, 0.4465)$ and standard 
deviation $\sigma = (0.2023, 0.1994, 0.2010)$ 
computed over the training set:

\begin{equation}
\hat{x} = \frac{x - \mu}{\sigma}
\label{eq:normalization}
\end{equation}

\subsubsection{Model Architecture}

A Convolutional Neural Network is trained from 
scratch on CIFAR-10 as the primary model for 
memorization analysis. The architecture consists 
of three convolutional blocks, each containing 
a convolutional layer, batch normalization, ReLU 
activation, and max pooling, followed by two 
fully connected layers for classification. The 
model is trained using the Adam optimizer defined 
in Equations \ref{eq:adammoments} and 
\ref{eq:adamupdate}, with a learning rate of 
$\eta = 0.001$, batch size of 64, and 
cross-entropy loss defined in Equation 
\ref{eq:crossentropy}.

\subsubsection{Memorization Analysis}

To identify memorized examples, the memorization 
score defined in Equation \ref{eq:memscore} is 
computed for each training example. Examples with 
a memorization score above a threshold $\tau$ are 
classified as memorized:

\begin{equation}
\mathcal{M} = \{(x_i, y_i) \in \mathcal{D} : 
\text{mem}(f, \mathcal{D}, i) \geq \tau\}
\label{eq:memset}
\end{equation}

where $\mathcal{M}$ denotes the set of memorized 
examples. The threshold $\tau$ is determined 
empirically based on the distribution of 
memorization scores across the training set.

\subsubsection{Adversarial Evaluation}

To assess the security implications of memorization, 
adversarial examples are generated using FGSM and 
PGD as defined in Equations \ref{eq:fgsm} and 
\ref{eq:pgd} respectively. The perturbation 
magnitude is set to $\epsilon = 0.03$ for both 
attacks, with PGD using a step size of 
$\alpha = 0.007$ and $T = 40$ iterations. 
Adversarial examples are generated for both 
memorized examples $\mathcal{M}$ and 
non-memorized examples to compare their 
relative vulnerability.

\subsubsection{Audit Pipeline}

The complete audit pipeline consists of the 
following sequential stages:

\begin{enumerate}
\item \textbf{Training:} Train the CNN on 
CIFAR-10 using the Adam optimizer and 
cross-entropy loss.
\item \textbf{Memorization Analysis:} Compute 
memorization scores for all training examples 
and identify the memorized set $\mathcal{M}$.
\item \textbf{Adversarial Evaluation:} Generate 
adversarial examples using FGSM and PGD for 
both memorized and non-memorized examples.
\item \textbf{XAI Attribution:} Apply GradCAM, 
Integrated Gradients, SHAP, and LIME to 
memorized examples and their adversarial 
counterparts.
\item \textbf{Analysis:} Compare attribution 
maps of memorized versus non-memorized examples 
to reveal instance-specific spurious features.
\end{enumerate}



